\chapter{Meneliti, Menerbitkan, dan Mematenkan}
\label[appendix]{lamp:riset}

Lampiran ini untuk pembaca yang ingin melangkah dari membaca ke meneliti. Isinya bersandar pada kuliah Patent and Literature
Searching yang saya ikuti pada musim dingin 2024/25, di semester pertama. Kuliahnya membahas dua jalan yang dipakai ilmu dan teknologi untuk membagikan temuan: publikasi ilmiah dan paten. Keduanya punya
aturan main yang berbeda, dan peneliti AI hampir pasti bersentuhan dengan keduanya.

\section{Publikasi ilmiah}

Kuliah membuka dengan kalimat yang jujur: publikasi adalah mata uang komunitas akademik. Peneliti dinilai dari makalah yang ia
terbitkan di jurnal bereputasi, dan dari seberapa sering makalah itu dikutip. Publikasi juga punya fungsi yang lebih mulia:
melaporkan hasil, memicu riset lanjutan, menghemat sumber daya karena percobaan yang sama tidak perlu diulang di tempat lain, dan
memberi pegangan bagi pembuat kebijakan dan industri.

Makalah ilmiah, apa pun bidangnya, mencoba menjawab serangkaian pertanyaan yang sama. Apa yang saya kerjakan, secara ringkas? Apa
masalahnya? Bagaimana saya mempelajari atau menyelesaikannya? Apa yang saya temukan? Apa artinya? Apa yang saya pelajari dan apa
yang sebaiknya dikerjakan berikutnya? Siapa yang membantu dan siapa yang membiayai? Karya siapa yang saya rujuk? Susunan abstrak,
pendahuluan, metode, hasil, diskusi, dan kesimpulan adalah cara baku untuk menjawabnya, dan setiap jurnal menetapkan detailnya dalam
pedoman penulis.

Sebelum terbit, makalah melewati \istilah{peer review}: editor mengirimkannya ke beberapa ahli di bidang yang sama, yang menilai apakah
metodenya benar, hasilnya didukung bukti, dan sumbangannya cukup baru. Di bidang machine learning, konferensi seperti NeurIPS, ICML,
dan ICLR memainkan peran yang di bidang lain dipegang jurnal, dan banyak makalah beredar lebih dulu sebagai pracetak di arXiv sebelum
atau tanpa melewati review. Pracetak mempercepat pertukaran gagasan, tetapi pembaca harus ingat bahwa belum ada orang lain yang
memeriksanya.

Kuliah juga mencatat apa yang \emph{tidak} biasanya diterbitkan: hasil negatif, percobaan yang gagal mereproduksi hasil yang sudah
terbit, dan temuan yang tampak bertentangan dengan konsep yang mapan. Akibatnya, literatur yang terbit cenderung lebih optimistis
dari kenyataan, sebuah bias publikasi. Pelajarannya bagi peneliti machine learning: satu makalah yang melaporkan metode baru
mengungguli pembanding belum berarti metode itu selalu lebih baik, dan reproduksi independen berharga.

\subsection{Mengukur dampak}

Ukuran dampak yang paling dikenal adalah \istilah{h-index}: nilai $h$ terbesar sehingga seorang peneliti punya $h$ makalah yang
masing-masing dikutip paling sedikit $h$ kali. Contoh di kuliah adalah daftar makalah dengan jumlah kutipan 25, 20, 10, 8, 7, 6, lalu
makalah-makalah yang dikutip lebih sedikit. Enam makalah teratas masing-masing dikutip paling sedikit enam kali, tetapi tidak ada
tujuh makalah yang masing-masing dikutip paling sedikit tujuh kali, sehingga $h=6$. Ukuran ini menggabungkan produktivitas dan dampak,
tetapi buta terhadap urutan penulis, bergantung pada bidang, dan bisa dinaikkan dengan saling mengutip. \emph{Impact factor}
sebuah jurnal, rata-rata kutipan yang diterima makalah-makalahnya dalam dua tahun sebelumnya, punya masalah serupa. Angka-angka ini
berguna sebagai petunjuk kasar, dan menyesatkan kalau dijadikan tujuan, hukum Goodhart yang sama dengan reward hacking di
\cref{ch:keselamatan}.

\subsection{Menemukan dan mengakses literatur}

Basis data seperti Web of Science, Scopus, SciFinder untuk kimia, dan Google Scholar membantu menemukan makalah yang relevan, melacak
siapa yang mengutip siapa, dan menemukan makalah-makalah baru dari satu makalah lama yang penting. Banyak makalah dilindungi hak cipta
dan berbayar. Kuliah membahas jalan-jalan yang sah untuk mendapatkannya: langganan perpustakaan universitas, meminta salinan langsung
kepada penulis, yang hampir selalu senang mengirimkannya, dan versi \emph{open access}. Pada \emph{gold open access}, penerbit membuat
makalah gratis dibaca dan biaya ditanggung penulis atau lembaganya. Pada \emph{green open access}, penulis mengarsipkan sendiri
salinannya, misalnya di repositori universitas atau arXiv. Pada \emph{diamond open access}, baik penulis maupun pembaca tidak membayar,
karena biaya ditanggung lembaga atau komunitas.

Daftar pustaka buku ini dibangun dengan alat-alat yang sama. Setiap artikel diambil metadatanya langsung dari Crossref berdasarkan
DOI-nya atau dari arXiv berdasarkan nomornya, lalu dicocokkan dengan kata kunci judul yang diharapkan, sehingga judul, penulis, dan
tahun yang tercetak memang berasal dari catatan resmi, bukan dari ingatan. Buku diperiksa terhadap katalog perpustakaan terbuka atau
halaman hak ciptanya. Kebiasaan sederhana ini mencegah kesalahan yang umum, yaitu mengutip sesuatu yang tidak pernah ada atau
mengutip dengan detail yang salah, kesalahan yang semakin sering terjadi sejak orang meminta model bahasa menuliskan daftar pustaka.

\section{Paten}

\subsection{Apa itu paten}

Paten adalah salah satu hak kekayaan intelektual, bersama merek, desain industri, hak cipta, dan model utilitas. Menurut definisi
yang dipakai kuliah, paten adalah hak yang diberikan negara kepada penemu untuk melarang orang lain membuat, memakai, menjual, atau
mengimpor penemuannya di negara itu selama waktu terbatas, \emph{sebagai ganti} pengungkapan penemuan itu kepada publik. Pertukaran
inilah intinya. Penemu mendapat monopoli sementara, masyarakat mendapat pengetahuan yang lengkap, yang bebas dipakai siapa pun setelah
patennya berakhir.

Gagasannya tua. Kuliah menyebut kota Yunani Sybaris sekitar 500 sebelum Masehi, yang konon melindungi resep masakan baru selama
setahun, monopoli sepuluh tahun bagi penemu cara menenun sutra di Venesia abad ke-12, undang-undang paten Venesia tahun 1474 yang sudah
memuat prinsip-prinsip modern, yaitu baru, berguna, eksklusif, dan terbatas waktunya, Statute of Monopolies di Inggris tahun 1623, dan
Konvensi Paris tahun 1883 yang menjadikan sistem paten internasional.

\subsection{Syarat dan siklus hidup}

Penemuan hanya bisa dipatenkan kalau memenuhi tiga syarat. Pertama, \emph{baru secara mutlak}: belum menjadi bagian dari keadaan
teknik sebelumnya di mana pun di dunia, dalam bentuk apa pun, termasuk makalah, presentasi konferensi, tesis, atau unggahan di
internet. Penemuan yang pernah diterbitkan di Jepang tidak lagi baru di Eropa, walau hampir tidak ada orang Eropa yang membacanya.
Kedua, \emph{langkah inventif}: tidak boleh jelas bagi seorang ahli di bidangnya. Ketiga, \emph{dapat diterapkan di industri}. Selain
itu, penemuan harus diungkapkan cukup jelas sehingga ahli di bidangnya bisa melaksanakannya.

Syarat kebaruan punya akibat yang sering mengejutkan peneliti. Di Eropa, publikasi oleh penemunya sendiri juga merusak kebaruan.
Peneliti yang mempresentasikan temuannya di konferensi minggu ini, lalu mengajukan paten bulan depan, sudah kehilangan haknya. Urutannya
harus dibalik: ajukan dulu, terbitkan kemudian. Kuliah menyebut jalan pintas yang murah untuk itu: kantor paten Austria, dan beberapa
kantor lain, menerima permohonan paten sementara dengan biaya resmi sekitar 50 euro. Permohonan ini tidak diperiksa dan tidak
diterbitkan, tetapi mengamankan tanggal prioritas, dan dalam setahun bisa ditingkatkan menjadi permohonan yang lengkap.

Setelah diajukan, permohonan diterbitkan sekitar 18 bulan setelah tanggal prioritas, diperiksa oleh kantor paten, dan kalau diberikan,
berlaku paling lama 20 tahun sejak tanggal pengajuan, asalkan biaya tahunan dibayar. Dalam tahun prioritas, permohonan yang sama bisa
diajukan di negara-negara lain dengan tetap memakai tanggal prioritas yang pertama.

\subsection{Membaca paten}

Bagian terpenting sebuah paten adalah daftar klaim bernomor, karena klaimlah, bukan gambar atau uraian, yang menentukan apa yang
dilindungi. Klaim independen berdiri sendiri, klaim dependen merujuk ke klaim lain dan menambahkan pembatasan. Satu kata kecil
menentukan luasnya. Klaim yang memakai ``comprising'' bersifat terbuka: produk yang punya semua unsur yang disebut, ditambah unsur
lain, tetap termasuk. Klaim yang memakai ``consisting of'' bersifat tertutup: hanya unsur yang disebut, tidak lebih. Kuliah
menyarankan urutan membaca: halaman pertama untuk jenis dokumen, tanggal prioritas, dan abstrak, lalu gambar dan contoh kalau
bidangnya belum dikenal, baru kemudian klaim.

Paten juga diklasifikasikan menurut bidang teknik, misalnya dengan Cooperative Patent Classification dan International Patent
Classification, sehingga pencarian keadaan teknik sebelumnya bisa menyisir satu kelas sekaligus. Tugas kuliah meminta kami menulis
ulasan sebuah paten: data bibliografisnya, kelasnya, mengapa tanggal prioritasnya lebih awal dari tanggal pengajuannya, apa
penemuannya, dan apa yang sebenarnya dilindungi klaimnya. Latihan ini mengajarkan bahwa uraian yang panjang di sebuah paten sering
jauh lebih luas dari perlindungan yang diberikan klaim-klaimnya.

Ketika paten dilanggar, pemiliknya bisa menggugat. Kuliah membahas tiga cara menghitung ganti rugi: royalti yang wajar, yaitu nilai pasar
sebuah lisensi; keuntungan yang hilang, yang harus dibuktikan sebagai akibat langsung pelanggaran; dan keuntungan yang diperoleh
pelanggar. Ancaman gugatan inilah yang memberi paten ``gigi''-nya.

\subsection{Paten dan AI}

AI menimbulkan dua pertanyaan baru bagi sistem paten. Yang pertama adalah apakah penemuan berbasis perangkat lunak dan machine learning
bisa dipatenkan. Di Eropa, program komputer ``sebagai program'' dikecualikan, tetapi penemuan yang diimplementasikan dengan komputer bisa
dipatenkan kalau ia memberi sumbangan teknis, misalnya metode machine learning yang dipakai untuk mengendalikan proses industri atau
mengolah sinyal fisis dengan cara baru. Batasnya tidak selalu jelas, dan banyak keputusan bergantung pada rumusan klaim. Yang kedua
adalah apakah AI bisa menjadi penemu. Dalam kasus DABUS, seorang pemohon mencantumkan sistem AI sebagai penemu satu-satunya. Kantor
Paten Eropa menolaknya pada 2020, dewan bandingnya menguatkan penolakan itu pada 2021, dan Mahkamah Agung Inggris pada 20 Desember
2023 memutuskan dengan suara bulat bahwa penemu harus manusia. Pertanyaan tentang siapa pemilik penemuan yang sebagian besar
dihasilkan AI, dan bagaimana menilai langkah inventif ketika ``ahli di bidangnya'' sudah dibantu AI, masih terbuka.

\sumberbab{Lampiran ini mengikuti kuliah Patent and Literature Searching (hak kekayaan intelektual dan sejarah paten, syarat paten,
siklus hidup paten, membaca dan mencari paten, strategi klaim, sengketa paten, publikasi ilmiah dan evaluasinya, mencari literatur
ilmiah) beserta tugas ulasan paten. Fakta tentang kasus DABUS diperiksa dengan laporan putusan resminya.}
